\section{Data and selection}\label{sec:data}

Which flights are compared, and what is measured on them?

\paragraph{Flights.} The flights are the cleaned tracks of the FFVL archive, both disciplines (the tracks of the old
repository, on the disk of the lab). The flight table of the reference run of the group-identification note gives, for
each flight, the interval of its logger, its departure (the first fix farther than 3 km from the takeoff) and its
landing. It holds \UniFunnelParaRunTable{} paraglider and \UniFunnelHangRunTable{} hang glider flights. A flight is
\emph{eligible} when its logger interval is at most $h_{\max}=$ \UniHMax{} s (so that every instant of the grid below lies in
a continuous stretch), when it has positions on the grid and when it departs. No group label is used: every flight is
taken, whoever it flew with.

\paragraph{Declared open, and open.} The FFVL catalogue gives the type of each flight. A flight is
\emph{declared open} when the pilot declared a distance flight: free distance, or distance with 1, 2 or 3 turn points. Triangles,
quadrilaterals and out-and-returns are declared closed. The \emph{opening} of a flight is
\[
  o \;=\; \frac{d}{L},\qquad d=\lVert \mathbf r_{\rm landing}-\mathbf r_{\rm takeoff}\rVert,\quad L=\text{path flown},
\]
with $\mathbf r_{\rm takeoff}$ and $\mathbf r_{\rm landing}$ the first and last fix of the flight and $L$ the length of its
track (Fig.~\ref{fig:opening}). A flight is \emph{open} when it is declared open and $o>$ \UniMinOpening{} (strictly): this
drops the declared open flights that come back near their takeoff. Table~\ref{tab:funnel} gives the flights at each step.

\begin{figure}
\centering
\includegraphics[width=\textwidth]{opening_sketch.pdf}
\caption{A flight is open when the distance from its takeoff to its landing, $d$ (dashed), is more than a tenth of the path
flown $L$. Two illustrative tracks; the openings printed are computed on them.}\label{fig:opening}
\end{figure}

\begin{table}
\centering
\small
\renewcommand{\arraystretch}{1.15}\setlength{\tabcolsep}{4pt}
\begin{adjustbox}{max width=\textwidth}
\begin{tabular}{lrr}
\toprule
Step & Para & Hang \\
\midrule
all flights of the run table & 156\,391 & 6\,094 \\
logger interval $\le$ 10 s, grid positions, a departure & 154\,598 & 6\,039 \\
eligible and declared open & 66\,845 & 2\,178 \\
declared open and opening $>$ 0.1 & 55\,879 & 1\,417 \\
open with at least half of the points on the DEM & 52\,583 & 1\,285 \\
pilots of the open flights & 5\,540 & 191 \\
\bottomrule
\end{tabular}
\end{adjustbox}
\caption{Flights at each step of the selection. 28 pilots have open flights in both disciplines.}\label{tab:funnel}
\end{table}


\paragraph{Pieces.} On the grid of $\Delta=$ \UniGridS{} s the track of a flight is cut into \emph{pieces}: the stretches inside one
segment of the cleaned track (the segments end at the gaps of the logger) and between the departure and the landing of the
flight. Positions are interpolated linearly on the grid within a segment, never across a gap. A piece shorter than \UniShortest{} s
is dropped. The pieces are whole: they are not cut by any label.

\paragraph{MSD, $\alpha(\tau)$ and $H$.} Let a piece have positions $\mathbf r_0,\dots,\mathbf r_{n-1}$ on the grid and
$\tau=m\Delta$ a lag with $m<n$. The sum and the number of origins at the lag $\tau$ are
\[
  S_p(\tau)=\sum_{k=0}^{n-1-m}\lVert \mathbf r_{k+m}-\mathbf r_k\rVert^2,\qquad N_p(\tau)=n-m .
\]
The MSD of a flight $f$ pools the origins of its pieces $p\in f$, and the MSD of a population is the mean over its flights, each
with the same weight:
\[
  \MSD_f(\tau)=\frac{\sum_{p\in f}S_p(\tau)}{\sum_{p\in f}N_p(\tau)},\qquad
  \overline{\MSD}(\tau)=\frac1{|P|}\sum_{f\in P}\MSD_f(\tau).
\]
The lags are \UniNLags{} lags of the earlier chapter of the thesis, geometric from \UniFirstLag{} s to \UniTauFixed{} s and then with the same
step up to \UniMaxLag{} s. The local exponent is the slope of $\log\overline{\MSD}$ against $\log\tau$ within $\pm0.25$
decades of $\tau$ (the three nearest lags when fewer lie there); we write
\[
  \alpha(\tau)=\frac{d\log\overline{\MSD}}{d\log\tau}\Big|_{\pm0.25\ \rm dec}=2H_{\rm loc}(\tau).
\]
The exponent $H$ is half the slope of the least-squares line through $\log\overline{\MSD}(\tau)$ against $\log\tau$, equal weight
per lag, over the lags from $10$ s to the longest lag $T$ of the curve:
$\overline{\MSD}(\tau)\propto\tau^{2H}$ (Fig.~\ref{fig:msd}). A line at $\alpha=2H$ marks it on the plots of $\alpha(\tau)$.

\begin{figure}
\centering
\includegraphics[width=0.55\textwidth]{msd_raw.pdf}
\caption{Over 10 to 1000 s the MSD of both disciplines follows a power law of exponent $2H$ (dashed fits). Open flights of
the fixed population of Section~\ref{sec:raw}, all terrain; bands 90\%.}\label{fig:msd}
\end{figure}

\paragraph{Two designs.} A lag $\tau$ needs a piece of at least $\tau/0.8$ s, so that a lag is at most 80\% of the piece (the
rule of the earlier chapter). Longer lags therefore leave out the flights that have only short pieces, and two
designs differ in how they treat them (Fig.~\ref{fig:design}).
\begin{itemize}
\item \textbf{A, fixed population.} A curve up to the lag $T$ uses only the flights that have a piece of at least $T/0.8$ s,
for every lag of the curve, so that the set of flights does not change along the curve. We use two values of $T$: the same
$T=$ \UniTauFixed{} s for every cell compared (``fixed''), and the longest lag of the grid at which both populations
compared keep at least \UniMinFlights{} flights (``longest''), stated for each comparison.
\item \textbf{B, all flights.} Each lag $\tau$ uses every flight that has a piece of at least $\tau/0.8$ s: the number of flights
falls with the lag.
\end{itemize}

\begin{figure}
\centering
\includegraphics[width=\textwidth]{design_sketch.pdf}
\caption{The fixed population (A) keeps the same flights at every lag, the all-flights design (B) lets them fall with the lag.
Seven illustrative flights, with the duration of their longest piece; dashed lines at $\tau/0.8$ for the lags marked.}\label{fig:design}
\end{figure}

\paragraph{Bands.} Bands are 90\% bands of a two-way bootstrap that resamples the pilots and the site-days (flights whose
takeoffs lie within \UniSiteKm{} km on the same day), with \UniNBoot{} draws each, shared by all populations so that the draws of a
difference between disciplines are paired (Appendix~\ref{app:bootstrap}). The pilots are those of the catalogue;
\UniPilotsBoth{} pilots have open flights in both disciplines.

\keep{Open flights of both disciplines, whole pieces on a 10-s grid, one MSD per flight, equal weight per flight;
$\alpha(\tau)$ is the local slope of the MSD, $H$ the slope from $10$ s to $T$. Design A compares the same flights at every lag,
design B all the flights each lag can use.}
